\documentclass[11pt,a4paper]{article}
\usepackage[T1]{fontenc}
\usepackage[utf8]{inputenc}
\usepackage{lmodern,microtype,amsmath,amsthm,amssymb}
\usepackage[left=26mm,right=26mm,top=24mm,bottom=25mm,headheight=14pt]{geometry}
\usepackage{xcolor,fancyhdr}
\usepackage[hypertexnames=false]{hyperref}
\definecolor{ink}{HTML}{204659}
\hypersetup{colorlinks=true,linkcolor=ink,citecolor=ink,urlcolor=ink,
pdftitle={Real collisions of zeros under the Gaussian backward heat flow},pdfauthor={Research manuscript prepared with ChatGPT for Keio}}
\pagestyle{fancy}\fancyhf{}\lhead{\small\itshape Real collisions under the Gaussian heat flow}\rhead{\small\thepage}
\renewcommand{\headrulewidth}{0.25pt}
\newtheorem{theorem}{Theorem}[section]
\newtheorem{lemma}[theorem]{Lemma}
\newtheorem{proposition}[theorem]{Proposition}
\newtheorem{corollary}[theorem]{Corollary}
\numberwithin{equation}{section}
\renewcommand{\Re}{\operatorname{Re}}\renewcommand{\Im}{\operatorname{Im}}
\setlength{\emergencystretch}{1.6em}
\title{Real collisions of zeros under the Gaussian backward heat flow}
\author{Research manuscript prepared with ChatGPT for Keio}
\date{3 November 2026}
\begin{document}
\maketitle
\begin{abstract}
Csordas and Escassut asked whether strict positivity of the first
Laguerre expression at every positive time of the backward heat flow
forces a strip entire function of order less than two to have only
real zeros. We prove their conjecture. If an initial zero is nonreal,
a real multiple zero occurs at a finite point and at a positive time
no later than half the square of the initial strip half-width.
The time bound is sharp. The main issue is ruling out loss of
collisions at horizontal infinity. We count the upper zeros of
canonical approximants with a rational weight whose pair interactions
cancel in a companion sum of imaginary parts. Applying the balance at
two spatial scales controls the tails of the collision measures.
The proof also gives a multiplicity count, a version for the full
canonical strip class on its natural time interval, and a weak
Laguerre test that detects collisions of either parity. For the Riemann
heat flow the weak test characterizes the Newman constant. Its
nonnegativity at all positive times remains to be proved.
\end{abstract}
\section{A Gaussian converse and a collision count}
Let $D=d/dz$ and $L(g)=(g')^2-gg''$. The strip class $\mathcal S(A)$
of Csordas and Escassut consists of nonzero real entire functions
\begin{equation}\label{eq:canonical}
 f(z)=e^{-a z^2}g(z),\qquad
 g(z)=Cz^d e^{bz}\prod_{\rho\ne0}(1-z/\rho)e^{z/\rho},
\end{equation}
where $a\ge0$, $b,C\in\mathbb R$, $C\ne0$, $d$ is a nonnegative
integer, $|\Im\rho|\le A<\infty$, and $\sum_{\rho\ne0}|\rho|^{-2}<\infty$.
All zeros are counted with multiplicity. Within this class, $f$ belongs
to the Laguerre--P\'olya class precisely when its zeros are all real.
At a real multiple zero the Laguerre expression vanishes.

\begin{theorem}[Csordas--Escassut Gaussian conjecture]\label{thm:CE}
Let $f\in\mathcal S(A)$ have order less than two and put
$F(t,z)=e^{-tD^2}f(z)$. If $L(F(t,\cdot))(x)>0$ for every
$t>0$ and $x\in\mathbb R$, all zeros of $f$ are real.
\end{theorem}
The source is the conjecture accompanying equation (3.3) of
\cite[pp.~336--337]{CE}: for their Gaussian symbol $\varphi(z)=e^{-z^2}$
one sets $t=\mu^2$. Their Open Problem~3.3 asks for a classification of symbols of order below two.
The same paper's Conjecture~3.5 concerns the
cosine symbol. The Gaussian implication follows from the geometric
statement below.

\begin{theorem}[Finite real collisions]\label{thm:collision}
Suppose $f$ has the representation \eqref{eq:canonical}, and put $F(t,z)=e^{-tD^2}f(z)$. Take
$0\le t<1/(4a)$, with no upper restriction if $a=0$.
If $f$ has a nonreal zero, then for some real $x_*$ and some $t_*$,
\[
 0<t_*\le \frac{A^2}{2+4aA^2},
 \qquad F(t_*,x_*)=F_z(t_*,x_*)=0 .
\]
More precisely, if $M$ is the number of initial upper-half-plane zeros,
then in the extended nonnegative integers
\begin{equation}\label{eq:exact-count}
 M=\sum_{(t,x)}\left\lfloor\frac{m(t,x)}{2}\right\rfloor .
\end{equation}
The sum is over positive times in the indicated domain and real
multiple zeros, and $m(t,x)$ is their spatial multiplicity.
\end{theorem}

For $f(z)=e^{-az^2}(z^2+A^2)$ with $A>0$, the bound for the first
collision is attained. For polynomial initial data with $a=0$ the
count follows by tracking finitely many roots. Strip contraction
puts all zeros on the real line by time $A^2/2$
\cite{deBruijn,Cardon}. For infinite products the corresponding
finite-dimensional statement alone gives no control over where the
approximating polynomials collide. We obtain that control by
comparing weighted counts at two different scales. No separation,
density assumption, or numerical verification of individual zeros
enters the argument.

\section{Canonical truncation and the range of the weights}
\label{sec:approx}
First let $a=0$. Keep the nonzero zeros in increasing disks,
retaining conjugate pairs, and let $f_n=e^{b_nz}P_n(z)$ be the
resulting canonical truncations, with $b_n$ real and $P_n$ a real
polynomial. They converge to $f$ locally uniformly. For each $\delta>0$
there is $C_\delta$ independent of $n$ such that
\[
                  |f_n(z)|\le C_\delta\exp(\delta |z|^2).
\]
Indeed $\log|(1-u)e^u|\le C|u|^2$ globally. A finite initial group of
factors has logarithmic maximum $O(1+|z|+\log(2+|z|))$,
and the remaining reciprocal-square sum can be made arbitrarily small.
The bound permits canonical products of order two and minimal type.

We use the Gaussian integral
\begin{equation}\label{eq:heat-int}
 F_n(t,z)=\frac1{\sqrt\pi}\int_{\mathbb R}e^{-v^2}
                          f_n(z+2i\sqrt t\,v)\,dv
\end{equation}
and the same formula for $F$. On a compact set of complex $t$ and $z$
the preceding bound with sufficiently small $\delta$ supplies an
integrable majorant. The integral is even in $\sqrt t$, so is entire
in $t$, and Taylor expansion gives $F_t=-F_{zz}$ (subscripts in
this identity denote partial derivatives). On each $0\le t\le T$
the transforms and their spatial derivatives converge locally uniformly
in $z$, uniformly in $t$. Also
\begin{equation}\label{eq:heat-growth}
 \sup_{n,\,0\le t\le T}\log^+\max_{|z|\le R}|F_n(t,z)|=o(R^2).
\end{equation}
For this last estimate bound $|z+2i\sqrt t\,v|^2$ by
$2|z|^2+8T v^2$ and take $\delta$ arbitrarily small.
The transform is injective on this growth class. Following
\eqref{eq:heat-int} by the real Gaussian displacement of the same
variance amounts to averaging $f(z+2\sqrt t\,(u+iv))$ against
$e^{-u^2-v^2}/\pi$. Angular averaging gives $f(z)$, with Fubini
justified by the same minimal growth. Hence $F(t,\cdot)\not\equiv0$.

For a fixed approximant,
\[
 F_n(t,z)=e^{b_n z-tb_n^2}\,(e^{-tD^2}P_n)(z-2tb_n).
\]
The polynomial form of the Gaussian strip theorem gives
$|\Im z|\le\sqrt{(A^2-2t)_+}$ for its zeros \cite{deBruijn,Cardon}.
Hurwitz's theorem gives the same conclusion for $F$. Put
$\ell=A+2$, $w_0=i\ell$. On $[0,T]$, $F_n(t,w_0)$ are uniformly
bounded away from zero and $|F_n'/F_n|(t,w_0)$ are uniformly bounded.
This follows by local convergence, compactness and nonvanishing at
$w_0$ for each of the finitely many discarded initial approximants.
Primes on spatial slices henceforth denote spatial derivatives.

Write $z_k=x_k+i y_k$ for the zeros of the current approximant at
time $t$. For $L\ge4\ell$ put $d_k=(L^2+x_k^2)^{1/2}$,
$h_k=L^2/d_k^2$. An important uniform bound is
\begin{equation}\label{eq:eps}
   \epsilon_L:=\sup_{n,\,0\le t\le T}\sum_k d_k^{-2}\longrightarrow0 .
\end{equation}
Along the imaginary axis the real exponential multiplier has constant modulus.
Consequently
\[
 \log\frac{|F_n(t,iL)|}{|F_n(t,i\ell)|}
 =\frac12\sum_k\log\left(1+
   \frac{(L-\ell)(L+\ell-2y_k)}{x_k^2+(\ell-y_k)^2}\right).
\]
Each summand is at least $cL^2/d_k^2$ after allowing for the factor
$1/2$, by $\log(1+v)\ge v/(1+v)$ for $v\ge0$ and $|y_k|\le A$.
Now \eqref{eq:eps} follows from \eqref{eq:heat-growth} and the
uniform lower bound at $i\ell$.

Define
\[
 g_L(z)=\frac{L^2}{L^2+z^2},\qquad r_L(z)=\frac z{L^2+z^2},\qquad
 B=b_n-\sum_k r_L(z_k)=\Re\frac{F'_n(t,iL)}{F_n(t,iL)} .
\]
The latter identity follows by averaging the logarithmic derivatives at $iL$ and $-iL$ and using conjugation. Since $b_n$ is real,
$-\Im(F_n'/F_n)(t,w_0)\ge2\sum_k|w_0-z_k|^{-2}$.
Subtract the logarithmic derivatives at $iL$ and $w_0$ and apply
Cauchy--Schwarz, using $|iL-z_k|\asymp d_k$, to obtain
\[
 |B|/L\le C/L+C\sqrt{\epsilon_L},\qquad
 \delta_L:=L^{-2}+\epsilon_L+\sup_{n,t}|B|/L\longrightarrow0 .
\]
All ensuing bounds are uniform over truncations and $t\in[0,T]$.
Their constants may depend on $f,A,T$.
\section{Weighted zero counts and their variation}\label{sec:balance}
Fix $T>A^2/2$ and $0\le s<T$. For this section abbreviate $g_L,r_L$ by $g,r$.

Write $U=\{i:y_i>0\}$ and let $H$ be the primitive of $g$ in
$|\Im z|<L$ with $H(0)=0$. Define
\[
 N(t)=\sum_{i\in U}\Re g(z_i),\qquad
 Y(t)=\sum_{i\in U}\Im H(z_i).
\]
On the strip $|y|\le A$, at $z=x+iy$, with
$d=(L^2+x^2)^{1/2}$ and $h=L^2/d^2$, direct differentiation gives
\[
\begin{gathered}
 \Re g\asymp h,\quad |g|\le C h,\quad
 |g'|\le C h/d,\quad |\Im g|\le C |y|h/d,\\
 |r|\le C/d,\qquad |r'|\le C/d^2,\qquad |\Im r|\le C|y|/d^2 .
\end{gathered}
\]
For example $\Re(L^2+z^2)=d^2-y^2\ge3d^2/4$ and
$g'=-2zL^2/(L^2+z^2)^2$. The imaginary-part bounds follow by
integration from the real axis. Similarly
$\Im H(x+iy)=\int_0^y\Re g(x+iv)\,dv\asymp yh$ for $y>0$.
In particular $N\asymp\sum_U h_i$ and $0\le Y\le CA N$.

At a time of simple zeros, the zero evolution equation is
$z'_i=2b_n+2\sum_{k\ne i}(z_i-z_k)^{-1}$. For $i\in U$ put
\[
 S_i=\sum_{k\notin U}\frac1{|z_i-z_k|^2},\qquad
 E_i=\sum_{k\notin U}\left(\frac1{z_i-z_k}+r(z_k)\right).
\]
Adding $r(z_k)$ to the internal interactions as well leaves the
one-body term $B+r(z_i)$. The compensated external terms satisfy
\begin{align}
 |E_i|&\le C d_i\sqrt{\epsilon_L S_i},\label{heat:Eabs}\\
 \Im E_i
 &\le -y_i S_i-\frac1{4y_i}+CA\epsilon_L.\label{heat:Eim}
\end{align}
For \eqref{heat:Eabs}, the numerator of a compensated summand is
$L^2+z_i z_k$ and its denominator is
$(z_i-z_k)(L^2+z_k^2)$. The numerator has modulus at most
$C d_i d_k$, so Cauchy--Schwarz applies. For \eqref{heat:Eim}
the imaginary part of the uncompensated sum is
$-y_i S_i+\sum_{k\notin U} y_k/|z_i-z_k|^2$; its conjugate
partner alone contributes $-1/(4y_i)$ to the second term.

The internal pairs disappear from the derivative of $Y$ by the
exact algebraic identity
\[
 \frac{g(z)-g(w)}{z-w}+g(z)r(w)+g(w)r(z)=0.
\]
For the derivative of $N$ apply $\partial_z+\partial_w$ to this
identity. The residual contribution of an internal pair before
taking real parts and restoring the factor two is
$-g(z)r'(w)-g(w)r'(z)$. Its total absolute value over internal
pairs is at most $C\epsilon_L N$.

The one-body term in $N'$ has modulus at most
$C(|B|/L+L^{-2})N$. For $Y'$ its imaginary part is bounded above
by $CA(|B|/L+L^{-2})N$, using $B$ real and the bounds on
$\Im g,\Im r$. Combining these estimates and
\eqref{heat:Eabs} gives
\begin{equation}\label{heat:Ndiff}
 |N'|\le C\delta_L N+
 C\sqrt{\epsilon_L}\sum_{i\in U}h_i\sqrt{S_i}.
\end{equation}
In $Y'$, the contribution $\Im(g_iE_i)$ is
$\Re g_i\,\Im E_i+\Im g_i\,\Re E_i$. The latter is bounded by
$C y_i h_i\sqrt{\epsilon_L S_i}$. Young's inequality absorbs it
in half of the negative $y_i h_i S_i$ term in \eqref{heat:Eim},
at the cost $C y_i h_i\epsilon_L$. Since
$y_i S_i+1/y_i\ge2\sqrt{S_i}$, we obtain
\begin{equation}\label{heat:Ydiff}
 Y'\le -c\sum_{i\in U}h_i\sqrt{S_i}+CA\delta_L N
\end{equation}
with $c>0$ independent of $L,n,t$.

\subsection{Real collisions of the approximants}
For a fixed polynomial the exceptional times are finite. Its
discriminant as a polynomial in $t$ is nonzero. Indeed, at large $t$, rescaling the space variable by $\sqrt t$ and
removing the real shift gives the heat transform of the leading
monomial in the limit, a Hermite polynomial with simple roots.
At a real zero of multiplicity $m$ at time $t_*$, rescale
$z-x_*=\sqrt{|t-t_*|}\,\xi$ after removing the moving real
exponential shift. The leading rescaled polynomial is
$e^{-\operatorname{sgn}(t-t_*)D_\xi^2}\xi^m$ up to a nonzero
scalar. For $t>t_*$ its roots are all real and simple. For
$t<t_*$ they are purely imaginary and simple, with the single
zero at the origin present exactly for odd $m$. Rouché's theorem
and reality give the same real/nonreal counts for nearby times.
Consequently $Y$ is continuous and $N$ loses exactly
$\lfloor m/2\rfloor g(x_*)$ at a real collision. A nonreal
collision causes no jump. Between exceptional times
\eqref{heat:Ndiff}--\eqref{heat:Ydiff} hold. Integration to their
endpoints is legitimate since polynomial zero branches have
convergent local Puiseux expansions. $N(s)$ denotes its
right-continuous value even if $s$ is exceptional.

Below we write $N_R$ or $N_L$ when the scale matters and $h_L(x)=L^2/(L^2+x^2)$. Let $V_L(t)=\int_s^t |N'(\tau)|\,d\tau$, excluding the finitely
many jump times. Integrating \eqref{heat:Ydiff}, using $Y(t)\ge0$
and $Y(s)\le CA N(s)$, then using \eqref{heat:Ndiff}, yields
\[
 V_L(t)\le C_A\sqrt{\epsilon_L} N(s)
             +C_A\delta_L(1+\sqrt{\epsilon_L})\int_s^t N(\tau)\,d\tau .
\]
All jumps are downward, so $N(t)\le N(s)+V_L(t)$. Grönwall's
inequality on the fixed horizon implies
$\sup_{[s,T]}N(t)\le C_{A,T}N(s)$ for all sufficiently large $L$,
and then
\[
 V_L(T)\le\eta_L N(s),\qquad
 \eta_L=C_{A,T}(\sqrt{\epsilon_L}+\delta_L)(1+\sqrt{\epsilon_L})
 \longrightarrow0 .
\]
Let $\nu_n$ be the counting measure of collisions on $(s,T]$
projected to the real axis, with weight $\lfloor m/2\rfloor$
for each real zero of multiplicity $m$. Since $N(T)=0$ by strip
contraction, we have
\begin{equation}\label{heat:loss}
 (1-\eta_L)N(s)\le \int_{\mathbb R}\frac{L^2}{L^2+x^2}\,d\nu_n(x)
                      \le (1+\eta_L)N(s).
\end{equation}

\section{Passing the collision count to the entire function}
Write $M$ for the number of upper zeros of $f$ and choose $T>A^2/2$.
Use the balance \eqref{heat:loss} with $s=0$. Its right-hand initial
count takes each retained upper zero $z$ of $f$ with weight $\Re g_L(z)$.
Consequently at fixed $L$ it converges as $n\to\infty$ to
\[
              N_L^f(0)=\sum_{\Im z>0,\ f(z)=0}\Re g_L(z).
\]
This is finite by \eqref{eq:canonical} and the comparison
$\Re g_L(z)\asymp L^2/(L^2+(\Re z)^2)$.

Applying the estimate at a second scale gives tightness for the
weighted measures. For $R\ge L$ and $|x|\ge R$,
$h_L(x)/h_R(x)\le 2L^2/R^2$, where $h_R(x)=R^2/(R^2+x^2)$.
By \eqref{heat:loss} used at scale $R$,
\[
 \int_{|x|\ge R}h_L(x)\,d\nu_n(x)
 \le \frac{2L^2}{R^2}(1+\eta_R) N_R(0)
 \le C L^2\epsilon_R\ \longrightarrow\ 0
 \qquad(R\to\infty)
\]
uniformly in $n$ at each fixed large $L$. This step also covers
order two and minimal type.

We describe the local limit of $\nu_n$. If $F$ has a real multiple
zero $x_*$ of multiplicity $m$ at time $t_*$, the heat equation
implies that
\[
 u^{-m}F(t_*+u^2,x_*+u v)
       \longrightarrow
       \frac{\partial_z^mF(t_*,x_*)}{m!}\,e^{-D_v^2}v^m
       \quad(u\to0)
\]
locally uniformly in $v$. This quotient is analytic in $(u,v)$ at
$u=0$ by the joint Taylor expansion, since the coefficient of
$u^{2j+k}v^k$ involves the spatial derivative of order $2j+k$.
Simplicity of the Hermite zeros gives $m$ local branches analytic
in $u$. At real times just after $t_*$ all their zeros are real
by conjugation. Just before $t_*$ exactly $\lfloor m/2\rfloor$ are
in the upper half-plane, by setting $u=i\sqrt{t_*-t}$. For odd $m$
the single central branch stays real, since all the other branches
are accounted for in conjugate pairs. Real collision points in
a compact space--time rectangle are therefore isolated and finite:
an accumulation point would itself be a collision by continuity and
would contradict this local description. At times $t>A^2/2$ there
can be no real collision because there are only real zeros at all
nearby earlier times.

Enclose an interior real collision in a small space--time cylinder with
a spatial disk centered on the real axis, with root-free lateral
boundary and simple zeros on the two time faces. Such cylinders
can be taken arbitrarily small. Uniform convergence and Rouch\'e's
theorem preserve the counts of upper and real zeros of $F$ on both
faces for all large $n$; the disks about simple real zeros are taken
symmetric about the axis. The upper count for $F_n$ on such a
cylinder changes only at real collisions. Its total decrease is
therefore exactly $\lfloor m/2\rfloor$, even if the approximating
collision splits into several events. At an initial multiple real
zero use instead a half-cylinder starting at time zero. The
canonical truncations have precisely the same initial zero there
for large $n$. Its whole cluster starts real, so no approximant
can lose an upper zero inside that cylinder. Collisions at time zero
are thus excluded. On a compact set outside the chosen
neighborhoods, simultaneous convergence of $F_n$ and $F'_n$ rules
out further real collisions. These facts imply convergence against
continuous test functions of compact spatial support to the measure
$\nu$ which assigns $\lfloor m/2\rfloor$ at each real collision of
$F$ in $(0,T]$, projected onto the real axis. In particular local
masses cannot accumulate at either time endpoint.

The tail estimate passes to $\nu$ by exhaustion with compactly
supported functions. Hence $\int h_L\,d\nu_n\to\int h_L\,d\nu$.
We can now pass to the limit in \eqref{heat:loss} to obtain
\[
 (1-\eta_L)N_L^f(0)\le \int h_L\,d\nu\le(1+\eta_L)N_L^f(0).
\]
As $L\to\infty$ the integral increases to $\nu(\mathbb R)$.
For finite $M$, $N_L^f(0)\to M$ term by term. For infinite $M$
the sum tends to infinity by retaining any fixed finite subset of
upper zeros. This proves \eqref{eq:exact-count} for $a=0$.
In particular, $M>0$ forces an actual finite real collision at a
time $0<t_*\le A^2/2$. This proves Theorem~\ref{thm:CE}, since a function of order
less than two has no quadratic Gaussian factor in
\eqref{eq:canonical} and strict positivity forbids such a collision.

For $a>0$ put $d_t=1-4at$. Gaussian conjugation gives
\begin{equation}\label{eq:lens}
 F(t,z)=d_t^{-1/2} e^{-az^2/d_t}
          (e^{-(t/d_t)D^2}g)(z/d_t),\qquad 0\le t<1/(4a).
\end{equation}
For instance both sides are analytic near $t=0$, satisfy
$F_t=-F_{zz}$ by differentiation, and have the same initial value and hence the same time Taylor expansion. The identity extends throughout the indicated interval by analytic
continuation. The Gaussian integral for the
left side is locally convergent there by the minimal growth of $g$.
The time $t/d_t$ runs from zero to infinity. Real rescaling and a
nonzero Gaussian preserve the zeros and their multiplicities, and
$t/d_t=A^2/2$ corresponds to $t=A^2/(2+4aA^2)$.
This completes Theorem~\ref{thm:collision}. Sharpness follows
from $e^{-sD^2}(z^2+A^2)=z^2+A^2-2s$ and \eqref{eq:lens}.

\section{A weak test at rational times and positions}\label{sect:weak}
Set
\[
 \Delta_f(t,x)=L(F(t,\cdot))(x)\,L(F_z(t,\cdot))(x).
\]
At a real collision of order $m$ at $(s,x_0)$, write
$u=\partial_z^mF(s,x_0)$. If $m=2r$, time Taylor expansion at $s-h$
for $h\downarrow0$ gives leading terms $u h^r/r!$ and
$u h^{r-1}/(r-1)!$ for $F$ and $F_{zz}$ at $x_0$, while $F_z=O(h^r)$.
The leading term of $L(F)$ is negative. The leading term of
$L(F_z)$ is the positive square of the leading term of $F_{zz}$.
For $m=2r+1$ apply the even-order calculation to $F_z$; the leading
term of $L(F)$ is then $u^2h^{2r}/(r!)^2$. Thus in both cases
\begin{equation}\label{eq:parity}
 \Delta_f(s-h,x_0)\sim
 \begin{cases}
 -u^4 h^{2m-3}/[r!(r-1)!^3],&m=2r,\\
 -u^4 h^{2m-3}/[r!^3(r-1)!],&m=2r+1.
 \end{cases}
\end{equation}
In particular every collision gives an \emph{open} set of negative
values immediately before its time. The first Laguerre form alone
can miss an odd collision: for $f(z)=z^3+z$ one has
$F(t,z)=z^3+(1-6t)z$ and $L(F)=3x^4+(1-6t)^2\ge0$ at all times.

\begin{corollary}\label{cor:weak}
For $f$ in \eqref{eq:canonical}, all its zeros are real if and only if
$\Delta_f(t,x)\ge0$ for all positive rational $t<1/(4a)$ and rational
$x$ (all positive rational $t$ if $a=0$).
\end{corollary}
\begin{proof}
A nonreal initial zero produces a collision by
Theorem~\ref{thm:collision}. Equation~\eqref{eq:parity} and continuity
give a rational witness. Conversely if $f$ has only real zeros, for
$a=0$ each heat-evolved truncation is a real linear exponential times
a real-rooted polynomial and so belongs to the Laguerre--P\'olya
class. Compact limits and differentiation preserve this class and
the first Laguerre inequality holds there. For $a>0$ the same conclusion
follows from \eqref{eq:lens}, since multiplying by a real Gaussian
and real rescaling also preserve the class.
\end{proof}

\section{What the Riemann specialization requires}
Let $H_t(z)=\int_0^\infty e^{t v^2}\Phi(v)\cos(zv)\,dv$ in the
standard normalization of \cite{RT,Poly}, so that
$H_0(z)=\Xi(z/2)/8$. The Newman constant $\Lambda$ is the threshold
for reality of all the zeros. Rodgers and Tao proved $\Lambda\ge0$
and Platt and Trudgian proved $\Lambda\le1/5$
\cite{RT,PT}. Write $\Delta(t,x)=L(H_t)(x)L(H_t')(x)$.
\begin{corollary}\label{cor:RH}
In this normalization,
\[
 \Lambda=\sup\bigl(\{0\}\cup
 \{t\in\mathbb Q:0<t<1/5,\ \Delta(t,x)<0
                        \text{ for some }x\in\mathbb Q\}\bigr).
\]
In particular RH is equivalent to nonnegativity in this rational
two-variable test.
\end{corollary}
\begin{proof}
For $t\ge\Lambda$, $H_t$ and $H_t'$ belong to the Laguerre--P\'olya
class and both factors are nonnegative. If $\Lambda>0$, take any
$0<a<\Lambda$. At time $a$ there is a nonreal zero. The function
$H_a$ has order one, is real on the real axis, and has bounded zero heights
\cite{Poly}. Theorem~\ref{thm:collision} applied to $H_a$, using
the heat semigroup, gives a real collision at some time $s>a$.
It satisfies $s\le\Lambda$, since a later collision would have
nonreal zeros immediately before it. Equation~\eqref{eq:parity}
gives an open set of negative witnesses in $a<t<s$, including
rational ones. Let $a\uparrow\Lambda$.
\end{proof}

For the Riemann family alone the existence of a last collision at
positive $\Lambda$ also follows from the uniform large-height
estimate in \cite[Theorem~1.5]{Poly}: zeros off the real axis have
$|\Re z|<\exp(C/t)$ for $0<t\le1/2$ and an absolute $C$.
A sequence of nonreal conjugate pairs as $t\uparrow\Lambda>0$
therefore has a finite limit subsequence, giving a multiple real
zero by Rouch\'e's theorem. Earlier work already uses exclusion of
positive-time collisions as an RH criterion \cite{Fathi}. At heat
time zero, the different condition using
$H_0(z+i\mu)+H_0(z-i\mu)$ for every $0<\mu<1$ is also an
established RH criterion \cite[Theorem~4.1, phase zero]{CSV}.
This relies on the known positive proportion of critical-line zeros,
so does not require a general density-free converse.

\paragraph{Unconditional information inside the flow.}
Equation~\eqref{eq:exact-count} implies that the number of nontrivial
zeta zeros with $\Im\rho>0$ and $\Re\rho>1/2$, with multiplicity,
equals
\[
       \sum_{0<t\le1/5,\ x>0}
                  \left\lfloor\frac{\operatorname{ord}_x H_t}{2}
                         \right\rfloor ,
\]
where the sum runs over real multiple zeros of $H_t$. Indeed the
change of variables is $\rho=1/2+iz/2$. The zeros of $H_0$ are
symmetric about both axes; it has no zeros on the imaginary axis,
and $H_t(0)>0$ by its integral. Half of each count therefore gives
the asserted symmetry reduction, also when both counts are infinite.
At each fixed positive time all but finitely many zeros of $H_t$
are real and simple \cite{KKL}. Applying the same count at that time
shows that only finitely many collision points can lie at later times.
Thus an infinite set of positive-time collisions would have times
tending to zero and spatial positions escaping every bounded interval.

We can also verify a region of the weak test without assuming RH.
For $0<t\le1/5$ put $X(t)=\exp(C/t)$ with the absolute constant from
\cite[Theorem~1.5]{Poly}. Then
\[
              \Delta(t,x)\ge0\quad\text{if } |x|\ge X(t)+2 .
\]
To check this, write $h=H_t$. Its nonreal zeros $c\pm id$ have
$|c|<X(t)$ and $0<d\le1$. Every nonreal zero of $h'$ is either a
zero of $h$ or lies in the closed disk $(x-c)^2+y^2\le d^2$ of
some such pair. Indeed, for $y>0$ each real zero gives a positive
term in $\partial_y\log|h(x+iy)|$, and each conjugate pair gives a
term with the sign of $(x-c)^2+y^2-d^2$, by logarithmic differentiation.
Outside all these disks the derivative is positive and $h'$ cannot
vanish. Thus the nonreal zeros of $h'$ have $|\Re z|\le X(t)+1$
and $|\Im z|\le1$. For either $k=h$ or $k=h'$ use its order-one
product on the real axis. A nonreal pair $c\pm id$ contributes
$2((x-c)^2-d^2)/((x-c)^2+d^2)^2$ to $L(k)/k^2$, and each real
zero contributes a positive reciprocal square. At a zero of $k$,
$L(k)=(k')^2$. Both forms are therefore nonnegative on the stated
range. This argument uses the high-height zero theorem and the
localization in \cite[Proposition~2.2]{Cs14}; it gives no numerical
value of $C$.

The interior region $0<t<1/5$, $|x|<X(t)+2$ still requires a sign proof.
Likewise the first-order test $L(H_0)\ge0$ on the entire real line
is the separate Open Problem~4.7 in \cite{Cs14}. The general collision
theorem establishes the Gaussian implication and the count. Neither
a proof of RH nor a numerical improvement of the Newman upper bound
is obtained here.

\begin{thebibliography}{9}\raggedright
\bibitem{CE} G.~Csordas and A.~Escassut,
The Laguerre inequality and the distribution of zeros of entire functions,
\emph{Ann. Math. Blaise Pascal} \textbf{12} (2005), 331--345.
\href{https://doi.org/10.5802/ambp.210}{doi:10.5802/ambp.210}.
\bibitem{deBruijn} N.~G.~de Bruijn, The roots of trigonometric integrals,
\emph{Duke Math. J.} \textbf{17} (1950), 197--226.
\href{https://doi.org/10.1215/S0012-7094-50-01720-0}{doi:10.1215/S0012-7094-50-01720-0}.
\bibitem{Cardon} D.~A.~Cardon, Complex zero strip decreasing operators,
\emph{J. Math. Anal. Appl.} (2015).
\href{https://doi.org/10.1016/j.jmaa.2015.01.026}{doi:10.1016/j.jmaa.2015.01.026}.
\bibitem{CSV} G.~Csordas, W.~Smith and R.~S.~Varga,
Level sets of real entire functions and the Laguerre inequalities,
\emph{Analysis} \textbf{12} (1992), 377--402.
\href{https://doi.org/10.1524/anly.1992.12.34.377}{doi:10.1524/anly.1992.12.34.377}.
\bibitem{Cs14} G.~Csordas, Fourier transforms of positive definite kernels
and the Riemann $\xi$-function, 2014 revision.
\href{https://arxiv.org/abs/1309.0055}{arXiv:1309.0055}.
\bibitem{Poly} D.~H.~J.~Polymath, Effective approximation of heat flow
evolution of the Riemann $\xi$ function, and a new upper bound for the
de Bruijn--Newman constant, \emph{Res. Math. Sci.} \textbf{6} (2019), article 31.
\href{https://doi.org/10.1007/s40687-019-0193-1}{doi:10.1007/s40687-019-0193-1}.
\bibitem{RT} B.~Rodgers and T.~Tao, The de Bruijn--Newman constant is
non-negative, \emph{Forum Math. Pi} \textbf{8} (2020).
\href{https://doi.org/10.1017/fmp.2020.6}{doi:10.1017/fmp.2020.6}.
\bibitem{PT} D.~Platt and T.~Trudgian, The Riemann hypothesis is true up to
$3\cdot10^{12}$, \emph{Bull. Lond. Math. Soc.} \textbf{53} (2021), 792--797.
\href{https://doi.org/10.1112/blms.12460}{doi:10.1112/blms.12460}.
\bibitem{KKL} H.~Ki, Y.-O.~Kim and J.~Lee, On the de Bruijn--Newman
constant, \emph{Adv. Math.} \textbf{222} (2009), 281--306.
\href{https://doi.org/10.1016/j.aim.2009.04.003}{doi:10.1016/j.aim.2009.04.003}.
\bibitem{Fathi} K.~Fathi, \emph{Zeros of the de Bruijn--Newman Flow},
October 2026 preprint package, Proposition~3.2 of the companion paper.
\href{https://doi.org/10.5281/zenodo.23604866}{doi:10.5281/zenodo.23604866}.
\end{thebibliography}
\end{document}
